% =====================================================================
\section{Topic 6: Directional Derivatives and the Gradient}\label{ch:grad}
% =====================================================================

\subsection*{Core formulas}

\begin{center}
\renewcommand{\arraystretch}{1.45}
\begin{tabularx}{\linewidth}{L{0.28\linewidth}Y}
\toprule
\textbf{Concept} & \textbf{Formula / fact}\\
\midrule
\fact{Partial derivatives}{$f_x$ = rate of change in the $\vi$ direction, $f_y$ = rate of change in the $\vj$ direction.}
\fact{Directional derivative (definition)}{$\Du f(x_0,y_0)=\displaystyle\lim_{h\to0}\frac{f(x_0+ha,\,y_0+hb)-f(x_0,y_0)}{h}$,
  \ $\vu=\ip{a,b}$ a \textbf{unit} vector.}
\fact{Formula}{$\Du f(x,y)=f_x(x,y)\,a+f_y(x,y)\,b$ \quad ($f$ differentiable).}
\fact{Angle form}{If $\vu$ makes angle $\theta$ with the positive $x$-axis, $\vu=\ip{\cos\theta,\sin\theta}$ and
  $\Du f=f_x\cos\theta+f_y\sin\theta$.}
\fact{Gradient}{$\grad f=\ip{f_x,f_y}=f_x\vi+f_y\vj$; \quad in 3D $\grad f=\ip{f_x,f_y,f_z}$.}
\fact{Dot-product form}{$\Du f=\grad f\cdot\vu=\norm{\grad f}\cos\theta$ ($\theta$ = angle between $\grad f$ and $\vu$).}
\fact{Maximum rate}{$\max\Du f=\norm{\grad f}$, in the direction $\vu=\dfrac{\grad f}{\norm{\grad f}}$ (fastest increase).}
\fact{Minimum rate}{$\min\Du f=-\norm{\grad f}$, in the direction $-\grad f$ (fastest decrease).}
\fact{Zero rate}{$\Du f=0$ when $\vu\perp\grad f$: the direction \emph{along the level curve}.}
\fact{Possible values}{Every value of $\Du f$ lies in $\bigl[-\norm{\grad f},\ \norm{\grad f}\bigr]$.}
\fact{Direction from $P$ toward $Q$}{$\vu=\dfrac{\vecPQ}{\norm{\vecPQ}}$ (always normalise!).}
\fact{Geometry}{$\grad f(P)$ is \emph{perpendicular} to the level curve $f=f(P)$ through $P$ and points toward
  \emph{larger} values of $f$.}
\fact{Tangent plane to $F(x,y,z)=k$}{$\grad F(P)\cdot\ip{x-x_0,y-y_0,z-z_0}=0$; normal line $\vr(t)=P+t\,\grad F(P)$.}
\fact{Chain rule (along a curve)}{$\dfrac{\dd}{\dd t}f(\vr(t))=\grad f(\vr(t))\cdot\vr'(t)$.}
\bottomrule
\end{tabularx}
\end{center}

\subsection*{Methods for the standard problems}

\begin{strategy}{Directional derivative at $P$ toward $Q$ (Quiz~2, Problem~2)}
\begin{enumerate}
\item $\grad f=\ip{f_x,f_y}$, then evaluate at $P$.
\item $\vv=\vecPQ=Q-P$, \ $\vu=\vv/\norm{\vv}$.
\item $\Du f(P)=\grad f(P)\cdot\vu$.
\item Fastest increase: unit vector $\grad f(P)/\norm{\grad f(P)}$. Maximum rate: $\norm{\grad f(P)}$.
\end{enumerate}
\end{strategy}

\begin{strategy}{Find the directions in which $\Du f(P)=c$ (Quiz~2, Problem~5)}
Let $\vu=\ip{a,b}$ and $\grad f(P)=\ip{p,q}$. Solve the system
\[ pa+qb=c,\qquad a^2+b^2=1. \]
Substitute the linear equation into the circle: 0, 1 or 2 unit vectors. Solutions exist only if
$\abs c\le\norm{\grad f(P)}$; if $\abs c=\norm{\grad f(P)}$ there is exactly one (namely
$\pm\grad f/\norm{\grad f}$). \ Shortcut: $\cos\theta=c/\norm{\grad f}$ gives the angle between
$\vu$ and $\grad f$; rotate $\grad f$ by $\pm\theta$.
\end{strategy}

\begin{strategy}{Points where the fastest increase is in a given direction $\vv$ (Quiz~2, Problem~6)}
Require $\grad f(x,y)=\lambda\vv$ with $\lambda>0$ (same direction, not opposite). E.g.\
$f=x^2+y^2-2x-4y$, $\vv=\vi+\vj$: $\ip{2x-2,2y-4}=\ip{\lambda,\lambda}$ gives
$2x-2=2y-4>0$, i.e.\ the points on the line $y=x+1$ with $x>1$.
\end{strategy}

\subsection*{Common mistakes}
\begin{itemize}
\item Using a direction vector that is not a unit vector. ($\ip{3,-4}$ must become $\ip{\frac35,-\frac45}$.)
\item Writing the \emph{maximum rate} as a vector. It is the number $\norm{\grad f}$; the
  \emph{direction} is the vector.
\item Evaluating the gradient at $Q$ instead of $P$.
\item Accepting $\lambda<0$ in ``direction of fastest increase is $\vv$'': that is fastest \emph{decrease}.
\end{itemize}
